\documentclass[conference]{IEEEtran}
\usepackage{graphicx}
\usepackage{amsmath}
\usepackage{amssymb}
\usepackage{booktabs}
\usepackage{microtype}
\usepackage{cite}
\usepackage{xcolor}

% Mark uncertain data in red
\newcommand{\tbd}[1]{\textcolor{red}{[TBD: #1]}}
\newcommand{\verify}[1]{\textcolor{orange}{#1 [VERIFY]}}

\title{A Cost-Aware Cascaded Pipeline for News Event Clustering with Selective LLM Verification}

\author{
    \IEEEauthorblockN{Jayadasan S, D. Menaga, Jason Peniel Raj S}
    \IEEEauthorblockA{St. Joseph's Institute of Technology, Chennai, India}
}

\begin{document}

\maketitle

\begin{abstract}
\textcolor{red}{[REWRITE AFTER GETTING REAL TEST RESULTS]}

News aggregators must identify when multiple sources report on the same event despite different wording.
Large language models (LLMs) can determine whether two headlines describe the same event, but exhaustive pairwise comparison is expensive at scale.
This paper presents NISE, a cascaded pipeline combining lexical filters with a five-signal evidence fusion algorithm (EFSA) to identify which pairs require LLM verification.
We evaluate on \tbd{N} test pairs from real RSS feeds.
The cascade achieves \tbd{accuracy}\% accuracy with \tbd{precision}\% precision and \tbd{recall}\% recall, reducing LLM calls by \tbd{X}\% compared to exhaustive verification.
A CPU-only Sentence-BERT alternative achieves \tbd{Y}\% accuracy with zero API cost on the same test set.
\end{abstract}

\begin{IEEEkeywords}
Cost-aware AI, event clustering, large language models, news aggregation.
\end{IEEEkeywords}

\section{Introduction}

\textcolor{red}{[HONEST FRAMING: Engineering study, no novel methods, empirical tradeoff analysis]}

When Apple announces earnings, 50+ outlets publish within an hour with different headlines.
News aggregators must group articles by event despite paraphrasing.
Lexical methods (Jaccard, n-grams) fail on synonym variation.
Dense embeddings (Sentence-BERT) capture semantics but require O(n²) comparisons.
LLMs achieve high accuracy but impose API costs.

This paper evaluates a \textbf{cascaded filtering strategy}: cheap lexical gates reject obvious non-matches, then an evidence fusion scorer ranks candidates, and only borderline cases reach LLM verification.

\textbf{Contribution:} An empirical analysis of cost-accuracy tradeoffs in production news clustering.
We report honest test-set performance including overfitting effects, compare with SBERT baseline on identical data, and measure actual LLM call reduction from stage logging.

\textbf{Key findings:}
\begin{itemize}
    \item \tbd{Insert honest test results here}
    \item \tbd{Conservative tuning led to overfitting - test recall dropped to X\%}
    \item \tbd{SBERT baseline outperforms cascade on test accuracy (Y\% vs Z\%)}
\end{itemize}

\section{Related Work}

\textcolor{red}{[KEEP EXISTING - IT'S MOSTLY OK]}

\subsection{Lexical Clustering}
Early text clustering used TF-IDF~\cite{salton1988} and bag-of-words.
Carta et al.~\cite{carta2022} applied inverted-index similarity for financial news, achieving sub-second latency but failing on synonym variation.

\subsection{Dense Embeddings}
Reimers and Gurevych~\cite{reimers2019sentence} introduced Sentence-BERT for semantic comparison.
While more effective than lexical methods, SBERT requires model inference and O(n²) comparisons.

\subsection{LLM Verification}
Recent work shows strong LLM performance on semantic tasks, but few studies analyze cost-accuracy tradeoffs at production scale with non-zero API pricing.

\section{Methodology}

\subsection{Pipeline Overview}
NISE processes articles from 18 RSS feeds across 14 sectors.
Each article is compared against articles from the previous 48 hours.

\textbf{Stage 1: Lexical Gates.}
Unigram Jaccard (τ=0.12) OR character 3-gram cosine (τ=0.25).
Pairs failing both are rejected.

\textbf{Stage 2: EFSA Scoring.}
Five-signal fusion: keyword IoU, 3-gram cosine, entity overlap, temporal decay, sector match.
\begin{equation}
S = 0.25 S_{key} + 0.30 S_{3g} + 0.25 S_{ent} + 0.10 S_{temp} + 0.10 S_{sec}
\end{equation}
Pairs with S ≥ 0.22 proceed to LLM.

\textbf{Stage 3: LLM Verification.}
Zero-shot pairwise comparison via Groq Llama 3.1-8B-Instant.

\subsection{SBERT Alternative}
CPU-only path using ONNX all-MiniLM-L6-v2 (threshold 0.50) for deployments without API access.

\section{Evaluation}

\subsection{Test Set}
\tbd{N=198 or N=1,022? Clarify which is test-only vs full}

\textcolor{red}{[BE EXPLICIT: Test set = 198 held-out pairs, Full = 1,022 including training]}

Benchmark constructed from \tbd{time period}, containing \tbd{pos} positive and \tbd{neg} negative pairs.
\tbd{Describe hard negatives: same sector, overlapping entities, temporal proximity}

Dual annotation: Cohen's κ = \verify{0.8454 ± 0.0402}

\subsection{Test Set Results (Primary)}

\textcolor{red}{[REPORT TEST SET AS PRIMARY RESULTS]}

Table~\ref{tab:test} shows performance on \tbd{N} held-out test pairs.

\begin{table}[h]
\centering
\caption{Performance on \tbd{N}-Pair Test Set}
\label{tab:test}
\begin{tabular}{lcccc}
\toprule
Method & Acc. & Prec. & Rec. & F1 \\
\midrule
Lexical Jaccard & \tbd{} & \tbd{} & \tbd{} & \tbd{} \\
Char 3-Gram & \tbd{} & \tbd{} & \tbd{} & \tbd{} \\
SBERT (τ=0.50) & \tbd{} & \tbd{} & \tbd{} & \tbd{} \\
\textbf{NISE Cascade} & \tbd{} & \tbd{} & \tbd{} & \tbd{} \\
SBERT+LLM & \tbd{} & \tbd{} & \tbd{} & \tbd{} \\
LLM-Only & \tbd{} & \tbd{} & \tbd{} & \tbd{} \\
\bottomrule
\end{tabular}
\end{table}

\textbf{Stage Pass Counts:}
\begin{itemize}
    \item Stage 1 passed: \tbd{X} pairs (\tbd{Y}\%)
    \item Stage 2 exceeded threshold: \tbd{Z} pairs (\tbd{W}\%)
    \item LLM calls made: \tbd{K} (\tbd{P}\% reduction vs exhaustive)
\end{itemize}

\subsection{Full Benchmark (In-Sample)}

\textcolor{red}{[LABEL AS IN-SAMPLE, NOT PRIMARY RESULTS]}

Table~\ref{tab:full} shows performance on full 1,022-pair benchmark (including training data).

\textcolor{orange}{WARNING: These results include training data and show signs of overfitting.}

\begin{table}[h]
\centering
\caption{In-Sample Performance (N=1,022, includes training)}
\label{tab:full}
\begin{tabular}{lcccc}
\toprule
Method & Acc. & Prec. & Rec. & F1 \\
\midrule
NISE Cascade & \verify{93.93\%} & \verify{100\%} & \verify{87.89\%} & \verify{93.56\%} \\
\bottomrule
\end{tabular}
\end{table}

\subsection{Cost Analysis}

\textcolor{red}{[FIX COST CALCULATION]}

API cost per call:
\begin{itemize}
    \item Input: 42 tokens × \$0.05/1M = \$0.0000021
    \item Output: 18 tokens × \$0.08/1M = \$0.0000014
    \item Total: \$0.00000354 per call
\end{itemize}

Cost per 1M comparisons:
\begin{itemize}
    \item LLM-only: 1M × \$0.00000354 = \$3.54
    \item NISE: \tbd{X} calls × \$0.00000354 = \tbd{\$Y}
    \item Savings: \$\tbd{Z} (\tbd{P}\% reduction)
\end{itemize}

\textcolor{red}{[REFRAME: Focus on latency and rate limits, not dollars]}

\textbf{More important than cost:}
\begin{itemize}
    \item Latency: \tbd{X}ms vs \tbd{Y}ms exhaustive
    \item Rate limits: \tbd{P}\% fewer API calls
    \item Scalability: Can process \tbd{X} articles/day
\end{itemize}

\subsection{Threshold Sweep}

\textcolor{red}{[ADD THIS TABLE - MISSING FROM ALL VERSIONS]}

\begin{table}[h]
\centering
\caption{EFSA Threshold Sweep on Test Set}
\label{tab:sweep}
\begin{tabular}{lccccc}
\toprule
τ & LLM Calls & Acc. & Prec. & Rec. & F1 \\
\midrule
0.18 & \tbd{} & \tbd{} & \tbd{} & \tbd{} & \tbd{} \\
0.20 & \tbd{} & \tbd{} & \tbd{} & \tbd{} & \tbd{} \\
0.22 & \tbd{} & \tbd{} & \tbd{} & \tbd{} & \tbd{} \\
0.25 & \tbd{} & \tbd{} & \tbd{} & \tbd{} & \tbd{} \\
\bottomrule
\end{tabular}
\end{table}

\subsection{Ablation Study}

\textcolor{red}{[ADD THIS TABLE - REVIEWERS EXPLICITLY ASKED]}

\begin{table}[h]
\centering
\caption{EFSA Leave-One-Out Ablation on Test Set}
\label{tab:ablation}
\begin{tabular}{lcccc}
\toprule
Configuration & Acc. & Rec. & F1 & Δ \\
\midrule
Full EFSA (5 signals) & \tbd{} & \tbd{} & \tbd{} & -- \\
w/o Keyword & \tbd{} & \tbd{} & \tbd{} & \tbd{} \\
w/o Headline 3-gram & \tbd{} & \tbd{} & \tbd{} & \tbd{} \\
w/o Entity & \tbd{} & \tbd{} & \tbd{} & \tbd{} \\
w/o Temporal & \tbd{} & \tbd{} & \tbd{} & \tbd{} \\
w/o Sector & \tbd{} & \tbd{} & \tbd{} & \tbd{} \\
\bottomrule
\end{tabular}
\end{table}

\subsection{Limitations}

\textcolor{red}{[BE HONEST]}

\begin{enumerate}
    \item \textbf{Overfitting:} \tbd{Large gap between in-sample (93.93\%) and test (X\%) suggests overfitting to training distribution.}
    \item \textbf{Low test recall:} \tbd{Test recall of X\% means lexical gates miss Y\% of true matches.}
    \item \textbf{SBERT outperforms:} SBERT standalone achieves \tbd{Z}\% vs cascade \tbd{W}\% on test set.
    \item \textbf{Single LLM:} Only evaluated Llama 3.1-8B; other models may differ.
    \item \textbf{English-only:} Generalization to other languages unknown.
\end{enumerate}

\section{Conclusion}

\textcolor{red}{[HONEST CONCLUSION]}

We evaluated a cascaded filtering strategy for news event clustering.
On a \tbd{N}-pair test set, the cascade achieves \tbd{X}\% accuracy with \tbd{Y}\% recall, reducing LLM calls by \tbd{Z}\%.

\textbf{Key findings:}
\begin{itemize}
    \item Lexical pre-filtering reduces API calls but at cost of recall
    \item \tbd{Conservative tuning led to overfitting (test recall X\% vs in-sample Y\%)}
    \item SBERT baseline outperforms cascade on test accuracy
    \item For deployments prioritizing precision, cascade achieves design goal
    \item For deployments prioritizing recall, SBERT is more effective
\end{itemize}

\textcolor{red}{[NO FAKE ACCEPTANCE SCORES, NO FABRICATED DEADLINES]}

\begin{thebibliography}{99}
\footnotesize

\bibitem{salton1988}
G. Salton and C. Buckley, ``Term-weighting approaches in automatic text retrieval,'' \textit{Inf. Process. Manage.}, vol. 24, no. 5, pp. 513--523, 1988.

\bibitem{carta2022}
S. Carta et al., ``Event detection and clustering on multi-source financial data streams,'' \textit{IEEE Trans. Knowl. Data Eng.}, vol. 34, no. 6, pp. 2845--2859, 2022.
\textcolor{orange}{[VERIFY: Check this exists on IEEE Xplore]}

\bibitem{reimers2019sentence}
N. Reimers and I. Gurevych, ``Sentence-BERT: Sentence embeddings using Siamese BERT-networks,'' \textit{Proc. EMNLP}, 2019, pp. 3982--3992.

\textcolor{red}{[VERIFY ALL OTHER REFERENCES BEFORE SUBMISSION]}

\end{thebibliography}

\end{document}
